% Lineare Gleichungen · Aufstellen · Prüfungs-Fokus MSA – bank/lineare-gleichungen/e4.jsonl (zusammenbau v0.6, Prüfungs-Fokus)
\einheitenkopf[e4]{Einheit 4 · Gleichungen aufstellen}
\zweigzeile{ab Kl. 6, je nach Buch bis Kl. 8 (am Gymnasium ab Kl. 7) · P10 ×7}
\setcounter{aufgabe}{0}

% Anlauf (ohne Prüfkennung)
\begin{aufgabe}{Aufstellen – Anlauf}
\begin{teile}
\teil „Eine Zahl plus 9 ergibt 23." – welche Gleichung passt? ($x$: die gesuchte Zahl) \\ \kreuz{$9x = 23$} \\ \kreuz{$x + 9 = 23$} \\ \kreuz{$x - 9 = 23$}
\teil Ich denke mir eine Zahl. Addiere ich 18, erhalte ich 45. Gleichung und Zahl? ($x$: die gesuchte Zahl) x = \leerfeld
\end{teile}
\end{aufgabe}

\begin{aufgabe}{Aufstellen – Prüfungsaufgaben}
\begin{teile}
\steil Ein Freibad hatte am Samstag 380 Gäste, der Eintritt kostet 5 € pro Person. Der Betreiber sagt: „Mit $x$ Gästen mehr hätten wir 2\,400 € eingenommen." Stelle eine Gleichung auf und berechne, wie viele Gäste mehr es hätten sein müssen. \hfill (P20*) \\ x = \leerfeld
\steil Ein Kino hat am Freitag 240 Karten zu je 9 € verkauft. Die Leiterin sagt: „Mit $x$ Karten mehr hätten wir 2\,700 € eingenommen." Stelle eine Gleichung auf und berechne, wie viele Karten mehr es hätten sein müssen. \hfill (P20*) \\ x = \leerfeld
\teil 500 Jugendliche wurden nach ihrem Lieblingssport gefragt. Tennis wurde viermal so oft genannt wie Volleyball. Berechne die Anzahlen für Tennis und Volleyball. \hfill (P18) \\

\sachtabelle{lc}{Sport & Anzahl}{Fußball & 140\\ Schwimmen & 90\\ Basketball & 70\\ Sonstiges & 40\\ Tennis & \leerzelle\\ Volleyball & \leerzelle}
Volleyball: \leerfeld , Tennis: \leerfeld
\teil 800 Gäste eines Festes wurden nach ihrem Lieblingsgetränk gefragt. Kaffee wurde fünfmal so oft genannt wie Kakao. Berechne die Anzahlen für Kaffee und Kakao. \hfill (P18) \\

\sachtabelle{lc}{Getränk & Anzahl}{Wasser & 150\\ Saft & 210\\ Tee & 90\\ Limo & 110\\ Kaffee & \leerzelle\\ Kakao & \leerzelle}
Kakao: \leerfeld , Kaffee: \leerfeld
\teil „Das Neunfache einer Zahl vermindert um 15 ist gleich 57." Gib die Gleichung an. \hfill (P16) ($x$: die gesuchte Zahl)
\teil „Das Sechsfache einer Zahl vermindert um 20 ist gleich 34." Gib die Gleichung an. \hfill (P16) ($x$: die gesuchte Zahl)
\teil „Das Dreifache einer Zahl vermindert um 7 ist gleich 29." Welche Gleichung passt? \hfill (P21) ($x$: die gesuchte Zahl) \\ \kreuz{$3x = -7 + 29$} \\ \kreuz{$3 - 7x = 29$} \\ \kreuz{$3x - 7 = 29$} \\ \kreuz{$29 - 7 = 3x$}
\teil „Das Vierfache einer Zahl vermindert um 9 ist gleich 23." Welche Gleichung passt? \hfill (P21) ($x$: die gesuchte Zahl) \\ \kreuz{$4 - 9x = 23$} \\ \kreuz{$4x - 9 = 23$} \\ \kreuz{$23 - 9 = 4x$} \\ \kreuz{$4x = -9 + 23$}
\end{teile}
\end{aufgabe}

\medskip
\uebersichtskasten{Gleichung aus Text: Unbekannte festlegen (x = …), Text in Term übersetzen, Gleichung lösen, Antwort mit Einheit. \\ „Das Dreifache einer Zahl, vermindert um 4, ist 11": 3x − 4 = 11 \\ Formeln umstellen: wie eine Gleichung – nach der gesuchten Größe auflösen. \quad U = 2 · (a + b) → a = U : 2 − b}
